% Part: sets-functions-relations
% Chapter: functions
% Section: functions-relations

\documentclass[../../../include/open-logic-section]{subfiles}


\begin{document}

\olfileid[fa-IR]{sfr}{fun}{rel}

\olsection{توابع به‌منزلهٔ روابط}

\begin{explain}
تابعی که !!{element}sی~$A$ را به !!{element}sی~$B$
نگاشت می‌کند، آشکارا رابطه‌ای میان $A$ و~$B$ تعریف می‌کند:
این رابطه میان $x$ و $y$ برقرار است اگر و تنها اگر $f(x) = y$.
حتی می‌توانیم---برای نمونه، اگر بخواهیم ریاضیات را بر مفاهیم بنیادیِ
کمتری بنا کنیم---تابع~$f$ را با این رابطه، یعنی با مجموعه‌ای از
زوج‌ها، \emph{یکی در نظر بگیریم}. آنگاه این پرسش پیش می‌آید:
کدام رابطه‌ها به این شیوه تابع تعریف می‌کنند؟
\end{explain}

\begin{defn}[نمودار یک تابع] فرض کنید $f\colon A \to B$ یک تابع باشد.
\emph{نمودارِ}~$f$ رابطهٔ $R_f \subseteq A \times B$ است که
به‌صورت زیر تعریف می‌شود:
\[
R_f = \Setabs{\tuple{x,y}}{f(x) = y}.
\]
\end{defn}

\begin{explain}
نمودار یک تابع بنا بر اصل گسترش به‌طور یکتا تعیین می‌شود.
از اصل گسترش برای مجموعه‌ها، اصل گسترش برای توابع نیز نتیجه می‌شود؛
همان اصلی که تا اینجا به‌طور ضمنی به کار برده‌ایم:
اگر $f$ و~$g$ دامنه و هم‌دامنهٔ یکسانی داشته باشند و برای هر
ورودی مقدار یکسانی بدهند، دو تابع برابرند.

به همین ترتیب، رابطه‌ای که هر مؤلفهٔ نخست را حداکثر به یک مؤلفهٔ دوم مربوط کند («تابع‌وار» باشد)، نمودار تابعی بر دامنهٔ همان رابطه است.
\end{explain}

\begin{prop}\ollabel{prop:graph-function}
فرض کنید $R \subseteq A \times B$ چنان باشد که:
\begin{enumerate}
\item اگر $Rxy$ و $Rxz$، آنگاه $y = z$؛ و
\item برای هر $x \in A$ یک $y \in B$ وجود دارد چنان‌که $\tuple{x,
y} \in R$.
\end{enumerate}
آنگاه $R$ نمودار تابع $f\colon A \to B$ است که به‌وسیلهٔ
$f(x) = y$ اگر و تنها اگر $Rxy$ تعریف می‌شود.
\end{prop}

\begin{proof}
به‌ازای هر ورودی در دامنه، شرط دوم وجود یک $y$ را تضمین می‌کند چنان‌که $Rxy$. اگر $z \neq
y$ دیگری وجود داشت چنان‌که $Rxz$، شرط مربوط به~$R$ نقض می‌شد. پس اگر
یک $y$ وجود داشته باشد چنان‌که $Rxy$، این $y$ یکتا است و بنابراین $f$
خوش‌تعریف است. آشکارا $R_f = R$.
\end{proof}

\begin{explain}
هر تابع $f\colon A \to B$ نموداری دارد؛ یعنی رابطه‌ای به‌صورت زیرمجموعه‌ای از $A
\times B$ که با $f(x) = y$ تعریف می‌شود. از سوی دیگر، هر رابطهٔ~$R
\subseteq A \times B$ که ویژگی‌های داده‌شده در
\olref{prop:graph-function} را داشته باشد، نمودار تابعی~$f \colon A \to
B$ است. به‌سبب این پیوند نزدیک میان توابع و
نمودارهایشان، می‌توانیم تابع را همان نمودار آن در نظر بگیریم. به
بیان دیگر، می‌توان توابع را با روابطی معین، یعنی با
مجموعه‌هایی معین از چندتایی‌ها، یکی در نظر گرفت. \oliflabeldef{sfr:rel:ref:sec}{بااین‌حال،
منظور از این «یکی در نظر گرفتن» همان است که در
\olref[sfr][rel][ref]{sec} توضیح دادیم: ادعا نمی‌کنیم که ماهیت توابع چیست؛
فقط می‌گوییم که برای کار ریاضی مفید است توابع را مجموعه‌هایی معین
\emph{در نظر بگیریم}. یکی از فایده‌های این روش آن است
که اکنون می‌توانیم}{اکنون می‌توانیم} همان‌گونه که روی روابط عملیات انجام دادیم،
روی توابع نیز عملیات مشابهی انجام دهیم (نگاه کنید به
\olref[sfr][rel][ops]{sec}). به‌ویژه:
\end{explain}

\begin{defn}\ollabel{defn:funimage}
فرض کنید $f \colon A \to B$ تابعی باشد و $C\subseteq A$ برقرار باشد.

\emph{تحدیدِ}~$f$ به~$C$ تابع
~$\funrestrictionto{f}{C}\colon C \to B$ است که به‌صورت زیر تعریف می‌شود:
$(\funrestrictionto{f}{C})(x) = f(x)$ برای هر $x \in C$. به
بیان دیگر، $\funrestrictionto{f}{C} = \Setabs{\tuple{x, y} \in R_f}{x \in
C}$.

\emph{اعمالِ}~$f$ بر~$C$ برابر است با $\funimage{f}{C} =
\Setabs{f(x)}{x \in C}$. این را \emph{تصویرِ}~$C$
تحت~$f$ نیز می‌نامیم.
\end{defn}

\begin{explain}
از این تعریف‌ها نتیجه می‌شود که $\ran{f} =
\funimage{f}{\dom{f}}$ برای هر تابع~$f$ برقرار است.
\oliflabeldef{sfr:rel:ops:sec}{این مفاهیم را می‌توان با تعریف‌های
\olref[sfr][rel][ops]{sec} و در نظر گرفتن توابع به‌صورت روابط مقایسه کرد.
\emph{یادداشت ویراستاری:} تعریف تصویر با تعریف قبلی سازگار است؛
اما قرارداد تحدید در اینجا تفاوت دارد. در تحدید تابع، فقط ورودی‌ها
به زیرمجموعهٔ دامنه محدود می‌شوند و هم‌دامنه تغییر نمی‌کند؛
در تعریف پیشینِ تحدید رابطه، هر دو مؤلفه به مجموعهٔ تعیین‌شده محدود می‌شدند.
دو عمل دیگر---وارون‌ها و حاصل‌ضرب‌های نسبی---به
جزئیات اندکی بیشتر نیاز دارند. آن جزئیات را در
\olref[inv]{sec} و
\olref[cmp]{sec} ارائه خواهیم کرد.}{}
\end{explain}

\end{document}
